\documentclass[../../main.tex]{subfiles}

\begin{document}

\section{开映射定理与闭图像定理}\label{section:ch13-开映射闭图像}

本节我们用 Baire 范畴定理建立分析无穷维 Banach 空间之间的线性算子的两个实质性工具——开映射定理与闭图像定理. Baire 范畴定理被用于证明下面的定理.

\begin{theorem}\label{theorem:ch13-象闭判据}
令 $X$ 和 $Y$ 为 Banach 空间, 线性算子 $T:X\to Y$ 是连续的. 则 $T(X)$ 是 $Y$ 的闭子空间当且仅当存在常数 $M>0$, 对给定 $y\in T(X)$, 存在 $x\in X$ 使得
\begin{align}
T(x)=y\text{ 且 }\|x\|\leqslant M\|y\|.\label{align:ch13-象闭判据}
\end{align}
\end{theorem}

\begin{proof}
我们首先假定存在 $M>0$ 使得\eqref{align:ch13-象闭判据}成立. 令 $\{y_{n}\}$ 为 $T(X)$ 中收敛到 $y_{\star}\in Y$ 的序列. 我们必须证明 $y_{\star}$ 属于 $T(X)$. 若必要的话, 通过选取一个子序列, 我们可以假设对所有 $n\geqslant 2$,
\begin{align*}
\|y_{n}-y_{n-1}\|\leqslant 1/2^{n}.
\end{align*}
根据 $M$ 的选择, 对每个自然数 $n\geqslant 2$, 存在向量 $u_{n}\in X$ 使得
\begin{align*}
T(u_{n})=y_{n}-y_{n-1}\text{ 且 }\|u_{n}\|\leqslant M/2^{n}.
\end{align*}
因此, 对 $n\geqslant 2$, 若我们定义 $x_{n}=\sum_{j=2}^{n}u_{j}$, 则
\begin{align}
T(x_{n})=y_{n}-y_{1}\label{align:ch13-逼近和}
\end{align}
且对所有 $k\geqslant 1$,
\begin{align}
\|x_{n+k}-x_{n}\|\leqslant M\cdot\sum_{j=n}^{\infty}1/2^{j}.\label{align:ch13-Cauchy}
\end{align}
但 $X$ 是一个 Banach 空间, 因此 Cauchy 序列 $\{x_{n}\}$ 收敛到向量 $x_{\star}\in X$. 在\eqref{align:ch13-逼近和}中取 $n\to\infty$ 的极限且用 $T$ 的连续性推出 $y_{\star}=T(x_{\star})+y_{1}$. 由于 $y_{1}$ 属于 $T(X)$, $y_{\star}$ 也如此. 因此 $T(X)$ 是闭的.

为证明逆命题, 假设 $T(X)$ 是 $Y$ 的闭子空间. 为了记号上的方便, 假设 $Y=T(X)$. 令 $B_{X}$ 与 $B_{Y}$ 分别表示 $X$ 与 $Y$ 的单位开球. 由于 $T(X)=Y$,
\begin{align*}
Y=\bigcup_{n=1}^{\infty}n\cdot T(B_{X})=\bigcup_{n=1}^{\infty}n\cdot\overline{T(B_{X})}.
\end{align*}
Banach 空间 $Y$ 有非空内部, 因此我们从 Baire 范畴定理推出存在自然数 $n$, 使得闭集 $n\cdot\overline{T(B_{X})}$ 包含一个开球, 我们将它写为 $y_{0}+[r_{1}\cdot B_{Y}]$. 因此,
\begin{align*}
r_{1}B_{Y}\subseteq n\,\overline{T(B_{X})}-y_{0}\subseteq 2n\,\overline{T(B_{X})}.
\end{align*}
因此, 若我们设 $r=2n/r_{1}$, 由于 $\overline{T(B_{X})}$ 是闭的, 我们得到 $\overline{B_{Y}}\subseteq r\cdot\overline{T(B_{X})}$. 因此, 由于 $\overline{B_{Y}}$ 是 $Y$ 中的闭单位球, 对每个 $y\in Y$ 与 $\varepsilon>0$, 存在 $x\in X$ 使得
\begin{align}
\|y-T(x)\|<\varepsilon\text{ 且 }\|x\|\leqslant r\cdot\|y\|.\label{align:ch13-迭代逼近}
\end{align}
我们宣称\eqref{align:ch13-象闭判据}对 $M=2r$ 成立. 事实上, 令 $y_{\star}$ 属于 $Y$, $y_{\star}\neq 0$. 根据\eqref{align:ch13-迭代逼近}, 其中 $\varepsilon=1/2\cdot\|y_{\star}\|$ 而 $y=y_{\star}$, 存在向量 $u_{1}\in X$ 使得
\begin{align*}
\|y_{\star}-T(u_{1})\|<1/2\|y_{\star}\|\text{ 且 }\|u_{1}\|\leqslant r\cdot\|y_{\star}\|.
\end{align*}
现在再次用\eqref{align:ch13-迭代逼近}, 这次 $\varepsilon=1/2^{2}\cdot\|y_{\star}\|$ 而 $y=y_{\star}-T(u_{1})$. 存在 $X$ 中的向量 $u_{2}$ 使得
\begin{align*}
\|y_{\star}-T(u_{1})-T(u_{2})\|<1/2^{2}\cdot\|y_{\star}\|\text{ 且 }\|u_{2}\|\leqslant r/2\cdot\|y_{\star}\|.
\end{align*}
我们继续这个过程, 归纳地选取 $X$ 中的序列 $\{u_{k}\}$, 使得对每个 $k$,
\begin{align*}
\|y_{\star}-T(u_{1})-T(u_{2})-\cdots-T(u_{k})\|<1/2^{k}\cdot\|y_{\star}\|\text{ 且 }\|u_{k}\|\leqslant r/2^{k-1}\cdot\|y_{\star}\|.
\end{align*}
对每个自然数 $n$, 定义 $x_{n}=\sum_{k=1}^{n}u_{k}$. 则根据 $T$ 的线性, 对每个 $n$,
\begin{align*}
\|y_{\star}-T(x_{n})\|\leqslant 1/2^{n}\cdot\|y_{\star}\|.
\end{align*}
\begin{align*}
\|x_{n+k}-x_{n}\|\leqslant r\cdot\|y_{\star}\|\cdot\sum_{j=n}^{\infty}1/2^{j}\text{ 且 }\|x_{n}\|\leqslant 2r\cdot\|y_{\star}\|.
\end{align*}
根据假设, $X$ 是完备的. 因此 Cauchy 序列 $\{x_{n}\}$ 收敛到 $X$ 中的向量 $x_{\star}$. 由于 $T$ 是连续的且范数是连续的,
\begin{align*}
T(x_{\star})=y_{\star}\text{ 且 }\|x_{\star}\|\leqslant 2r\cdot\|y_{\star}\|.
\end{align*}
因此\eqref{align:ch13-象闭判据}对 $M=2r$ 成立. 证明完成了.
\end{proof}

从拓扑空间 $X$ 到拓扑空间 $Y$ 的映射 $f:X\to Y$ 称为开的, 若 $X$ 中的每个开集的象在拓扑空间 $f(X)$ 中是开的, 其中 $f(X)$ 具有承袭自 $Y$ 的子空间拓扑. 因此从 $X$ 到 $Y$ 的连续一一映射 $f$ 是开的, 当且仅当 $f$ 是 $X$ 与 $f(X)$ 之间的拓扑同胚.

\begin{theorem}[开映射定理]\label{theorem:ch13-开映射}
令 $X$ 与 $Y$ 为 Banach 空间而线性算子 $T:X\to Y$ 是连续的. 则它的象 $T(X)$ 是 $Y$ 的一个闭子空间当且仅当算子 $T$ 是开的.
\end{theorem}

\begin{proof}
前一个定理告诉我们, 证明 $T$ 是开的当且仅当存在常数 $M>0$ 使得\eqref{align:ch13-象闭判据}成立. 令 $B_{X}$ 与 $B_{Y}$ 分别表示 $X$ 与 $Y$ 的开单位球. 我们从 $T$ 与范数的齐次性推出\eqref{align:ch13-象闭判据}等价于包含关系:
\begin{align*}
\overline{B_{Y}}\cap T(X)\subseteq M\cdot T(\overline{B_{X}}).
\end{align*}
根据齐次性, 该包含关系等价于使得 $B_{Y}\cap T(X)\subseteq M'\cdot T(B_{X})$ 的常数 $M'$ 的存在性. 因此, 我们必须证明 $T$ 是开的, 当且仅当存在 $r>0$ 使得
\begin{align}
[r\cdot B_{Y}]\cap T(X)\subseteq T(B_{X}).\label{align:ch13-开球包含}
\end{align}
首先假设算子 $T$ 是开的. 则 $T(B_{X})\cap T(X)$ 是 $T(X)$ 的包含 $0$ 的开子集. 因此存在 $r>0$, 使得 $r\cdot B_{Y}\cap T(X)\subseteq T(B_{X})\cap T(X)\subseteq T(B_{X})$. 因此\eqref{align:ch13-开球包含}对 $r$ 的这个选取成立.

为证明反命题, 假设\eqref{align:ch13-开球包含}成立. 令 $\mathcal{O}$ 为 $X$ 的开子集而 $x_{0}$ 属于 $\mathcal{O}$. 我们必须证明 $T(x_{0})$ 是 $T(\mathcal{O})$ 的内点. 由于 $x_{0}$ 是 $\mathcal{O}$ 的内点, 存在 $R>0$ 使得 $x_{0}+R\cdot B_{X}\subseteq\mathcal{O}$. 我们从\eqref{align:ch13-开球包含}推出, $T(X)$ 中的以 $T(x_{0})$ 为中心、半径为 $r\cdot R$ 的开球包含于 $T(\mathcal{O})$. 因此 $T(x_{0})$ 是 $T(\mathcal{O})$ 的内点.
\end{proof}

\begin{corollary}\label{corollary:ch13-逆连续}
令 $X$ 与 $Y$ 为 Banach 空间而 $T\in\mathcal{L}(X,Y)$ 是一对一与映上的. 则 $T^{-1}$ 是连续的.
\end{corollary}

\begin{proof}
算子 $T^{-1}$ 是连续的当且仅当算子 $T$ 是开的.
\end{proof}

\begin{corollary}\label{corollary:ch13-范数等价}
令 $\|\cdot\|_{1}$ 与 $\|\cdot\|_{2}$ 为线性空间 $X$ 上的范数, 使得 $(X,\|\cdot\|_{1})$ 与 $(X,\|\cdot\|_{2})$ 都是 Banach 空间. 假定存在 $c\geqslant 0$, 使得
\begin{align*}
\text{在 }X\text{ 上},\quad\|\cdot\|_{2}\leqslant c\cdot\|\cdot\|_{1}.
\end{align*}
则这两个范数等价.
\end{corollary}

\begin{proof}
对所有 $x\in X$, 定义恒等算子 $\operatorname{Id}:X\to X$ 为 $\operatorname{Id}(x)=x$. 根据假设,
\begin{align*}
\operatorname{Id}:(X,\|\cdot\|_{1})\to(X,\|\cdot\|_{2})
\end{align*}
是 Banach 空间之间的有界算子, 因此是连续的算子, 当然它既是一对一又是映上的. 根据开映射定理, 恒等算子的逆 $\operatorname{Id}:(X,\|\cdot\|_{2})\to(X,\|\cdot\|_{1})$ 也是连续的, 即它是有界的: 存在 $M\geqslant 0$, 使得
\begin{align*}
\text{在 }X\text{ 上},\quad\|\cdot\|_{1}\leqslant M\cdot\|\cdot\|_{2}.
\end{align*}
因此这两个范数等价.
\end{proof}

\begin{definition}[闭算子]\label{definition:ch13-闭算子}
赋范线性空间 $X$ 与 $Y$ 之间的线性算子 $T:X\to Y$ 称为闭的, 只要 $\{x_{n}\}$ 是 $X$ 中的序列,
\begin{align*}
\text{若 }\{x_{n}\}\to x_{0}\text{ 且 }\{T(x_{n})\}\to y_{0},\quad\text{则 }T(x_{0})=y_{0}.
\end{align*}
映射 $T:X\to Y$ 的图像是集合 $\{(x,T(x))\in X\times Y\mid x\in X\}$. 因此算子是闭的当且仅当它的图像是乘积空间 $X\times Y$ 的闭子空间.
\end{definition}

\begin{theorem}[闭图像定理]\label{theorem:ch13-闭图像}
令 $T:X\to Y$ 为 Banach 空间 $X$ 与 $Y$ 之间的线性算子. 则 $T$ 是连续的当且仅当它是闭的.
\end{theorem}

\begin{proof}
显然 $T$ 是闭的, 若它是连续的. 为证明反命题, 假设 $T$ 是闭的. 引入 $X$ 上的新范数 $\|\cdot\|_{\star}$ 为
\begin{align*}
\text{对所有 }x\in X,\quad\|x\|_{\star}=\|x\|+\|T(x)\|.
\end{align*}
算子 $T$ 的闭性等价于赋范线性空间 $(X,\|\cdot\|_{\star})$ 的完备性. 另一方面, 我们显然有
\begin{align*}
\text{在 }X\text{ 上},\quad\|\cdot\|\leqslant\|\cdot\|_{\star}.
\end{align*}
由于 $(X,\|\cdot\|_{\star})$ 与 $(X,\|\cdot\|)$ 都是 Banach 空间, 从前一个推论得出存在 $c\geqslant 0$, 使得
\begin{align*}
\text{在 }X\text{ 上},\quad\|\cdot\|_{\star}\leqslant c\cdot\|\cdot\|.
\end{align*}
因此对所有 $x\in X$,
\begin{align*}
\|T(x)\|\leqslant\|x\|+\|T(x)\|\leqslant c\|x\|.
\end{align*}
因此 $T$ 是有界的从而是连续的.
\end{proof}

\textbf{注} 令 $X$ 与 $Y$ 是 Banach 空间而算子 $T:X\to Y$ 是线性的. 为证明 $T$ 的连续性, 必要的是证明: 若在 $X$ 中 $\{x_{n}\}\to x_{0}$, 则在 $Y$ 中 $\{T(x_{n})\}\to T(x_{0})$. 闭图像定理提供了这一准则的巨大的简化. 它告诉我们为证明 $T$ 的连续性, 仅须对 $X$ 中满足 $\{x_{n}\}\to x_{0}$ 的序列 $\{x_{n}\}$ 与 $Y$ 中的 Cauchy 序列 $\{T(x_{n})\}$, 检验在 $Y$ 中 $\{T(x_{n})\}\to T(x_{0})$. 这一简化的用途体现在即将介绍的定理 11 的证明中.

令 $V$ 为线性空间 $X$ 的线性子空间. 用 Zorn 引理的论证(见习题 35), 证明存在 $X$ 的子空间 $W$ 具有以下直和分解:
\begin{align}
X=V\oplus W.\label{align:ch13-直和分解}
\end{align}
我们称 $W$ 为 $V$ 在 $X$ 中的线性补. 若 $X$ 的子空间在 $X$ 中有有限维的线性补, 则它称为在 $X$ 中有有限余维. 对 $x\in X$ 与分解式\eqref{align:ch13-直和分解}, 令 $x=v+w$, 其中 $v\in V$, $w\in W$. 定义 $P(x)=v$. 我们把以下证明留作代数的练习. $P:X\to X$ 是线性的,
\begin{align}
\text{在 }X\text{ 上},\quad P^{2}=P,\quad P(X)=V\text{ 且 }(\operatorname{Id}-P)(X)=W.\label{align:ch13-投影性质}
\end{align}
我们称 $P$ 为 $X$ 沿 $W$ 映上 $V$ 的投影. 我们把证明若 $P:X\to X$ 是任何满足 $P^{2}=P$ 的线性算子, 则
\begin{align}
X=P(X)\oplus(\operatorname{Id}-P)(X)\label{align:ch13-投影分解}
\end{align}
留作第二个代数的练习.

我们因此称使得 $P^{2}=P$ 的线性算子 $P:X\to X$ 为投影. 若 $P$ 是投影, 则 $(\operatorname{Id}-P)^{2}=\operatorname{Id}-P$. 因此 $\operatorname{Id}-P$ 也是投影.

现在假设线性空间 $X$ 是赋范的. 使得\eqref{align:ch13-直和分解}成立的 $X$ 的闭子空间 $W$ 称为 $V$ 在 $X$ 中的闭线性补. 一般来说, 很难确定一个线性子空间是否有闭线性补. 下一章的推论 8 告诉我们赋范线性空间的每个有限维子空间有闭线性补. 第 16 章的定理 3 告诉我们 Hilbert 空间的每个闭子空间有闭线性补. 关于闭线性补的存在性, 从现在开始我们有用投影的连续性刻画的以下准则.

\begin{theorem}\label{theorem:ch13-闭线性补}
令 $V$ 为 Banach 空间 $X$ 的闭子空间. 则 $V$ 在 $X$ 中有闭线性补当且仅当存在将 $X$ 映上 $V$ 的连续投影.
\end{theorem}

\begin{proof}
假定 $P$ 是定义在 $X$ 上的连续投影, 令 $\{x_{n}\}$ 为 $P(X)$ 中收敛到 $x$ 的序列. 则 $x_{n}=P(x_{n})\to P(x)$, 因此 $x=P(x)$. 这表明 $P(X)$ 是闭的. 现在假设存在 $X$ 映上 $V$ 的连续投影 $P$. 存在直和分解 $X=V\oplus(\operatorname{Id}-P)(X)$. 由于 $\operatorname{Id}-P$ 是连续投影, $(\operatorname{Id}-P)(X)$ 是闭的. 为证明反命题, 假设存在 $X$ 的闭子空间 $W$ 使得存在直和分解式\eqref{align:ch13-直和分解}. 定义 $P$ 为 $X$ 沿着 $W$ 映上 $V$ 的投影. 我们宣称 $P$ 是连续的. 闭图像定理告诉我们, 为证明此断言仅须证明算子 $P$ 是闭的. 令 $\{x_{n}\}$ 为 $X$ 中的序列, 使得 $\{x_{n}\}\to x_{0}$ 且 $\{P(x_{n})\}\to y_{0}$. 由于 $\{P(x_{n})\}$ 是闭集 $V$ 中收敛于 $y_{0}$ 的序列, 向量 $y_{0}$ 属于 $V$. 由于 $\{(\operatorname{Id}-P)(x_{n})\}$ 是闭集 $W$ 中收敛于 $x_{0}-y_{0}$ 的序列, 向量 $x_{0}-y_{0}$ 属于 $W$. 因此 $P(y_{0})=y_{0}$ 且 $P(x_{0}-y_{0})=0$. 因此 $y_{0}=P(x_{0})$. 从而算子 $P$ 是闭的.
\end{proof}

根据定理 8 和它的推论——开映射定理, 给出准则确定何时连续线性算子的象是闭的是有趣的. 以下定理给出一个这样的准则.

\begin{theorem}\label{theorem:ch13-有限余维闭象}
令 $X$ 和 $Y$ 为 Banach 空间而线性算子 $T:X\to Y$ 是连续的. 若 $T(X)$ 在 $Y$ 中有闭线性补, 则 $T(X)$ 在 $Y$ 中是闭的. 特别地, 若 $T(X)$ 在 $Y$ 中有有限余维, 则 $T(X)$ 在 $Y$ 中是闭的.
\end{theorem}

\begin{proof}
令 $Y_{0}$ 为 $Y$ 的闭子空间, 使得
\begin{align}
T(X)\oplus Y_{0}=Y.\label{align:ch13-直和Y}
\end{align}
由于 $Y$ 是 Banach 空间, $Y_{0}$ 也是. 考虑 Banach 空间 $X\times Y_{0}$, 其中笛卡儿积上的线性结构依分量定义, 而范数定义为
\begin{align*}
\text{对所有 }(x,y)\in X\times Y_{0},\quad\|(x,y)\|=\|x\|+\|y\|.
\end{align*}
则 $X\times Y_{0}$ 是 Banach 空间. 定义线性算子 $S:X\times Y_{0}\to Y$ 为
\begin{align*}
\text{对所有 }(x,y)\in X\times Y_{0},\quad S(x,y)=T(x)+y.
\end{align*}
则 $S$ 是 Banach 空间 $X\times Y_{0}$ 映上 Banach 空间 $Y$ 的连续线性映射. 从定理 8 得出, 存在 $M>0$ 使得对每个 $y\in Y$ 存在 $(x,y')\in X\times Y_{0}$, 使得
\begin{align*}
T(x)+y'=y\text{ 且 }\|x\|+\|y'\|\leqslant M\cdot\|y\|.
\end{align*}
但由于 $T(X)\cap Y_{0}=\{0\}$, $y'=0$. 因此
\begin{align*}
T(x)=y\text{ 且 }\|x\|\leqslant M\cdot\|y\|.
\end{align*}
我们再次用定理 8 得出 $T(X)$ 是 $Y$ 的闭子空间. 最后, 由于赋范线性空间的每个有限维闭子空间是闭的, 若 $T(X)$ 有有限余维, 它是闭的.
\end{proof}

\textbf{注} 有限维赋范线性空间上的所有线性算子是连续的、开的, 且有闭的象. 本节的结果仅对定义在无穷维 Banach 空间上的线性算子有意义, 在这种情形下算子的连续性不蕴涵象是闭的. 我们把以下证明留作练习: 对 $\{x_{n}\}\in\ell^{2}$, 定义为
\begin{align*}
T(\{x_{n}\})=\{x_{n}/n\}
\end{align*}
的算子 $T:\ell^{2}\to\ell^{2}$ 是连续的, 但没有闭的象且不是开的.

\begin{exercise}\label{exercise:ch13-29}
令 $X$ 为有限维赋范线性空间而 $Y$ 是赋范线性空间. 证明每个线性算子 $T:X\to Y$ 是连续与开的.
\end{exercise}

\begin{exercise}\label{exercise:ch13-30}
令 $X$ 为 Banach 空间而 $P\in\mathcal{L}(X,X)$ 是投影. 证明 $P$ 是开的.
\end{exercise}

\begin{exercise}\label{exercise:ch13-31}
令 $T:X\to Y$ 是 Banach 空间 $X$ 与 $Y$ 之间的连续线性映射. 证明 $T$ 是开的当且仅当 $X$ 中的开单位球在 $T$ 下的象在 $Y$ 的原点的邻域是稠密的.
\end{exercise}

\begin{exercise}\label{exercise:ch13-32}
令 $\{u_{n}\}$ 为 Banach 空间 $X$ 中的序列. 假定 $\sum_{k=1}^{\infty}\|u_{k}\|<\infty$. 证明存在 $x\in X$ 使得
\begin{align*}
\lim_{n\to\infty}\sum_{k=1}^{n}u_{k}=x.
\end{align*}
\end{exercise}

\begin{exercise}\label{exercise:ch13-33}
令 $T$ 为赋范线性空间 $X$ 到无穷维赋范线性空间 $Y$ 的线性算子. 证明 $T$ 是连续的当且仅当 $\ker T$ 是 $X$ 的闭子空间.
\end{exercise}

\begin{exercise}\label{exercise:ch13-34}
假定 $X$ 是 Banach 空间, 算子 $T\in\mathcal{L}(X,X)$ 是开的而 $X_{0}$ 是 $X$ 的闭子空间. $T$ 在 $X_{0}$ 上的限制 $T_{0}$ 是连续的. $T_{0}$ 必须是开的吗?
\end{exercise}

\begin{exercise}\label{exercise:ch13-35}
令 $V$ 为线性空间 $X$ 的线性子空间. 用以下方法证明 $V$ 在 $X$ 中有线性补.
\begin{enumerate}[(i)]
\item 若 $\dim X<\infty$, 令 $\{e_{i}\}_{i=1}^{n}$ 为 $V$ 的基. 将 $V$ 的这个基延拓得到 $X$ 的基 $\{e_{i}\}_{i=1}^{n+k}$. 接着定义 $W=\operatorname{span}[\{e_{n+1},\cdots,e_{n+k}\}]$.
\item 若 $\dim X=\infty$, 将 Zorn 引理用于所有使得 $V\cap Z=\{0\}$ 的 $X$ 的子空间 $Z$ 的族 $\mathcal{F}$, 用集合的包含关系对该族排序.
\end{enumerate}
\end{exercise}

\begin{exercise}\label{exercise:ch13-36}
证明\eqref{align:ch13-投影性质}与\eqref{align:ch13-投影分解}.
\end{exercise}

\begin{exercise}\label{exercise:ch13-37}
令 $Y$ 为赋范线性空间. 证明 $Y$ 是 Banach 空间当且仅当存在 Banach 空间 $X$ 与将 $X$ 映上 $Y$ 的连续线性开映射.
\end{exercise}

\end{document}